% =====================================================================
%  PREAMBLE dung chung cho ca hai ban:
%    main_revision.tex   -> ban SACH (nop cho tap chi)
%    main_annotated.tex  -> ban CO DANH DAU (de thay va AE doi chieu)
%  Sua o day thi ca hai ban cung doi -- khong co ban nao lech ban nao.
% =====================================================================

% =====================================================================
%  TETC-2026-05-0252 -- BAN REVISION
%
%  Dung lai tu paper1.pdf vi file .tex goc khong con. Tieu de, danh sach
%  tac gia va don vi GIU NGUYEN: IEEE cam doi tac gia neu khong co dong y
%  bang van ban cua EiC, va reviewer trich dan theo so muc cu.
%
%  Bien dich:  pdflatex -> pdflatex   (bibliography la thebibliography,
%              khong can BibTeX)
%
%  Cac manh duoc \input o duoi deu co the kiem doc lap:
%      python runners/audit_c4.py        -> 100/100
%      python runners/audit_figures.py   ->  36/36
% =====================================================================

\usepackage{cite}
\usepackage{amsmath,amssymb,amsfonts}
\usepackage{amsthm}
\usepackage{algorithmic}
\usepackage{graphicx}
\usepackage{textcomp}
\usepackage{xcolor}
\usepackage{booktabs}
\usepackage{multirow}
\usepackage{array}
\usepackage{url}
\usepackage{bm}
\usepackage[hidelinks]{hyperref}

\graphicspath{{./}}

% =====================================================================
%  Tham so dat float. Mac dinh cua LaTeX chi cho float chiem <=70% dinh
%  trang va bat chua >=20% chu, nen hinh rong ca trang gan nhu luon bi day
%  di cho khac; day du xa thi no roi xuong cuoi bai va xen vao danh muc tai
%  lieu. Noi rong ra thi hinh dat duoc ngay trang no duoc nhac toi.
% =====================================================================
\setcounter{topnumber}{2}
\setcounter{bottomnumber}{1}
\setcounter{totalnumber}{3}
\setcounter{dbltopnumber}{2}
\renewcommand{\topfraction}{0.92}
\renewcommand{\bottomfraction}{0.5}
\renewcommand{\textfraction}{0.07}
\renewcommand{\floatpagefraction}{0.80}
\renewcommand{\dbltopfraction}{0.92}
\renewcommand{\dblfloatpagefraction}{0.80}

% Trang chi co float: mac dinh LaTeX can giua theo chieu doc, nen hai hinh
% bi keo xa nhau bang mot khoang trang lon. Ep can tren, khoang cach co
% dinh, phan du don xuong duoi.
\makeatletter
\setlength{\@fptop}{0pt}
\setlength{\@fpsep}{10pt plus 1fil}
\setlength{\@fpbot}{0pt plus 2fil}
\setlength{\@dblfptop}{0pt}
\setlength{\@dblfpsep}{10pt plus 1fil}
\setlength{\@dblfpbot}{0pt plus 2fil}
\makeatother

% Moi truong dinh ly. Ban da nop dung Proposition cho ca nhung su that hien
% nhien; ban nay chi con Lemma (dan xuat rieng cua bai) va Definition.
\newtheorem{definition}{Definition}
\newtheorem{lemma}{Lemma}
\newtheorem{assumption}{Assumption}
\newtheorem{problem}{Problem}
\theoremstyle{remark}
\newtheorem{remark}{Remark}

\newcommand{\ZZ}{\textsc{zz}}
\newcommand{\Zmap}{\textsc{z}}
\newcommand{\QSVM}{QSVM-\ZZ}
\newcommand{\Fmac}{F_1^{\mathrm{macro}}}

% =====================================================================
%  DANH DAU THAY DOI so voi ban da nop (TETC bat buoc nop ban danh dau).
%
%  Mot nguon, hai ban:
%    main_revision.tex   -> ban SACH, khong hien gi
%    main_annotated.tex  -> bat co \annottrue, hien nhan o le trai
%
%  Dung \revnote{loai}{ghi chu} ngay sau moi \subsection. Cac loai:
%    new      muc hoan toan moi, ban da nop khong co
%    rewrite  co trong ban cu nhung viet lai
%    fix      SUA LOI cua ban cu
%    keep     giu nguyen y cua ban cu
% =====================================================================
\newif\ifannot
\definecolor{tagnew}{HTML}{0B6E4F}
\definecolor{tagrew}{HTML}{1F3D6B}
\definecolor{tagfix}{HTML}{A3341F}
\definecolor{tagkeep}{HTML}{6B6B6B}
\newcommand{\revnote}[2]{%
  \ifannot
    \par\nobreak\vspace{1pt}%
    {\footnotesize\sffamily
     \csname revtag@#1\endcsname~\textcolor{tagkeep}{#2}}%
    \par\nobreak\vspace{2pt}%
  \fi}
\expandafter\def\csname revtag@new\endcsname{%
  \colorbox{tagnew}{\textcolor{white}{\bfseries MOI}}}
\expandafter\def\csname revtag@rewrite\endcsname{%
  \colorbox{tagrew}{\textcolor{white}{\bfseries VIET LAI}}}
\expandafter\def\csname revtag@fix\endcsname{%
  \colorbox{tagfix}{\textcolor{white}{\bfseries SUA LOI}}}
\expandafter\def\csname revtag@keep\endcsname{%
  \colorbox{tagkeep}{\textcolor{white}{\bfseries GIU}}}

% Ten cong khai de dung trong bang chu giai cua main_annotated.tex
\newcommand{\revtagnew}{\csname revtag@new\endcsname}
\newcommand{\revtagrewrite}{\csname revtag@rewrite\endcsname}
\newcommand{\revtagfix}{\csname revtag@fix\endcsname}
\newcommand{\revtagkeep}{\csname revtag@keep\endcsname}


